\section{Comparing different affine structures}
The printed proof of source Lemma~3.5 assumes that two invariant affine
structures are nested merely because their line sizes are ordered. The
following comparison does not make that assumption.

\begin{lemma}[An intrinsic translation group]\label{intrinsictranslations}
Let $G$ act primitively on a set $X$ of size $p^d$, $d\geq2$, and contain
a regular elementary abelian normal subgroup $V$. If $G$ preserves a
Desarguesian affine-space structure on $X$, then $V$ is its translation
group.
\end{lemma}
\begin{proof}
The point stabilizer acts irreducibly on $V$: a proper invariant subgroup
would give a nontrivial block system of its cosets. Thus $V$ is minimal
normal in $G$. In the given affine structure, project $G$ to its cyclic
field-automorphism group. The kernel of the restriction to $V$ is normal
in $G$ and is nontrivial, because an elementary abelian group of rank at
least two cannot inject into a cyclic group. Minimality puts all of $V$
in the affine linear subgroup.

Take linear parts there. A finite $p$-subgroup of a linear group in
characteristic $p$ has a nonzero fixed vector: orbit counting on the
vector space makes the number of fixed vectors divisible by $p$.
The common fixed space $W$ of the image of $V$ is invariant under the
entire semilinear image of $G$, since $V\lhd G$. If $W$ were proper,
its affine cosets would be a nontrivial $G$-invariant partition of $X$.
Thus $W$ is the whole vector space. The linear image of $V$ is trivial,
so $V$ consists of translations; equality follows from the common order
$p^d$.
\end{proof}

\begin{proposition}[Scalar rigidity]\label{scalarrigidity}
In Lemma~\ref{intrinsictranslations}, identify $X$ with $V$ by choosing
an origin. Suppose its stabilizer contains a perfect subgroup $N$ and
$\End_{\F_p N}(V)=\F_q$. Then every $G$-invariant Desarguesian affine
structure has line size at most $q$. Its scalar field is an actual
subfield of this commuting algebra. If the natural $\F_q$-structure is
$G$-invariant, it has largest line size among these structures.
\end{proposition}
\begin{proof}
All structures have the same translation group by the lemma, hence the
same addition on $V$. In another structure over $\F_Q$, the field action
of $N$ is a homomorphism to a cyclic group and is trivial because $N$ is
perfect. Thus $N$ is $\F_Q$-linear. Every scalar of $\F_Q$ commutes with
$N$, giving an embedding of that actual scalar field into
$\End_{\F_p N}(V)=\F_q$. In particular $Q\leq q$. This argument
constructs the common scalar setting rather than assuming a prior
partition of one pencil by the other.
\end{proof}

For natural special-linear modules of dimension at least three and natural
symplectic modules of dimension at least four in the thick range $q\geq3$,
the perfectness and the commuting-algebra calculations above apply.
For the even-field natural $G_2(q)$ module, the commuting algebra follows
from an order argument. Retain its standard faithful six-dimensional
representation, transitivity on nonzero vectors, simplicity for $q\geq4$,
and order $q^6(q^6-1)(q^2-1)$. Transitivity gives irreducibility over
$\F_p$. Its commuting algebra is a finite division ring, hence a field
containing $\F_q$, of relative degree dividing six. Relative degrees three
or six would embed the group into $\GL_2(q^3)$ or $\GL_1(q^6)$, impossible
by the $p$-parts of their orders. Relative degree two would embed this
perfect simple group into $\operatorname{SL}_3(q^2)$ and then, with
trivial central intersection, into $\operatorname{PSL}_3(q^2)$.
The ratio of the latter order to $q^6(q^6-1)(q^2-1)$ is
$(q^2+1)/3$, not an integer since $q$ is a power of two. Thus the
commuting algebra is exactly $\F_q$. The $q=2$ derived-group case
has been proved separately in Proposition~\ref{binaryg2}.

For the one-dimensional branch no comparison of affine structures is
needed to bound line size. A primitive pencil has prime size by
Lemma~\ref{prime}; if its line size is $p^f$ and point size $p^d$, that
lemma gives $d=u^{a+1}$ and $f=u^a$ for a prime $u$. Every nondegenerate
affine space on $p^d$ points has line size $p^b$ where $b$ is a proper
divisor of $d$, so $b\leq u^a=f$. This proves maximality even against
non-Desarguesian affine planes.

\begin{proposition}[The natural even-field exceptional pencil]\label{naturalg2}
Retain the classical six-dimensional symplectic representation of
$G_2(q)$, $q\geq4$ even, in which the points of the split Cayley hexagon
are exactly the vector lines, and their stabilizers are maximal
parabolic subgroups. Every group between this group and its full
$\F_2$-linear normalizer is primitive but not two-transitive on vector
lines. With translations its affine action is locally primitive but not
locally two-transitive. The analogous binary derived-group assertion
is Proposition~\ref{binaryg2}.
\end{proposition}
\begin{proof}
Transitivity and maximality of the stated point stabilizers imply
primitivity; these are the exact retained building inputs, attributed
to the classical hexagon identification used in source Proposition~3.6.
The commuting-field calculation above makes every normalizer element
semilinear over $\F_q$. Remove its field automorphism and pull back the
preserved nondegenerate alternating form. The resulting form is again
$G_2(q)$-invariant. Its defining operator relative to the original form
commutes with $G_2(q)$, so it is a scalar by the calculated commuting
algebra. Hence the entire normalizer preserves perpendicularity on
vector lines. There are distinct perpendicular and nonperpendicular
pairs in dimension six, so none of these groups is two-transitive on
the pencil. The retained transitivity of the natural group on nonzero
vectors implies transitivity of a line stabilizer on its nonzero
vectors. Translations then give two-transitivity on each affine line.
Proposition~\ref{translation} proves the remaining assertions.
\end{proof}

\begin{lemma}[Large vector-stabilizer orbits]\label{largeorbits}
Let $V=\F_q^n$, $q\geq3$. If $n\geq3$ and
$N=\operatorname{SL}_n(q)$, then $N_v$, for $0\ne v\in V$, is
transitive on $V\setminus\F_qv$, of size $q^n-q>q^{n/2}$.
If $n\geq4$ is even and $N=\operatorname{Sp}_n(q)$, every
$N_v$-orbit outside $\F_qv$ has size at least
$q^{n-1}-q>q^{n/2}$.
Consequently a $G$-invariant affine-space structure on $V$, where
$G_0\geq N$, has no line larger than $q$.
\end{lemma}
\begin{proof}
For the special-linear assertion, extend $(v,w)$ and $(v,w')$ to bases
for any $w,w'\notin\F_qv$. Map the bases to each other and correct
the determinant on a third basis vector. This fixes $v$ and sends $w$
to $w'$. For the symplectic assertion, symplectic basis extension gives
one orbit $v^\perp\setminus\F_qv$, of size $q^{n-1}-q$, and one
orbit $\{w:\langle v,w\rangle=c\}$, of size $q^{n-1}$, for each
$c\ne0$. The first transitivity follows by extending an ordered
independent isotropic pair; the second follows by extending the
hyperbolic pair $(v,c^{-1}w)$. These constructions account for every
vector outside $\F_qv$. The displayed strict inequalities hold already
at $n=3$ and $n=4$, respectively, and increase thereafter.

Here the comparison is only with affine spaces, not arbitrary linear
spaces. Every nondegenerate finite affine space on $|V|=q^n$ points
has line size at most $\sqrt{|V|}$: a plane has exactly $k^2$ points
if its line size is $k$, and higher dimension gives at least a plane.
For an alternative line $L$ through $0$ and $v$, uniqueness of the line
through these two points makes $L$ invariant under $N_v$.
It cannot contain any of the above outside orbits. Thus
$L\subseteq\F_qv$, and its size is at most $q$.
\end{proof}

\begin{theorem}[Maximal line size for a primitive affine pencil]
\label{affinemax}
Let a fixed faithful permutation group $G$ on a finite set $X$ act
flag-transitively on an affine space $\Gamma$ with at least two lines
through each point. If $G_x$ is primitive on its pencil, then
$\Gamma$ has largest line size among all $G$-invariant affine-space
structures on $X$ on which $G$ is flag-transitive.
This recovers source Lemma~3.5 by contraposition. Its proof retains the
precise BDDKLS--Huybrechts affine classification and the BDD exclusion
of local point-primitivity for the full L\"uneburg-plane group.
\end{theorem}
\begin{proof}
A plane already has line size $\sqrt{|X|}$, the universal affine bound
just proved. If the lines of $\Gamma$ have size two, a larger-line
structure would partition each pencil of the circle into the sets
$L'\setminus\{x\}$, contradicting Proposition~\ref{circles}.
We may therefore suppose $\Gamma=\operatorname{AG}(n,q)$,
$n\geq3$, $q\geq3$; the coordinatization of affine spaces of dimension
at least three is the classical retained input.

Use exactly the classification interface quoted in source Section~3.2.1:
a flag-transitive group is one-dimensional, belongs to the excluded
$\operatorname{AG}(4,3)$ last-derived-$2.A_5$ case, or is two-transitive
with one of Huybrechts' six stated natural representation types.
The excluded case is a retained impossibility, not inferred here from
its abstract group name. In the one-dimensional case the prime-pencil
and prime-power argument immediately preceding this theorem proves the
line-size bound, including comparison with non-Desarguesian planes.

In the two-transitive cases write $G=V\rtimes G_0$. The three small
exceptional representation types have imprimitive pencils by the
explicit exclusions in the classical-pencils section. For a natural
special-linear, symplectic or even-field $G_2$ normal subgroup over
$\F_{q^t}$, the commuting-algebra arguments make $G_0$ normalize that
actual scalar field. If $t>1$, its lines partition the $\F_q$-pencil
nontrivially, contrary to Lemma~\ref{field}. (The special-linear
one-dimensional case belongs to the one-dimensional alternative.)
Thus the remaining modules are natural over $\F_q$ itself. In the
special-linear and symplectic cases Lemma~\ref{largeorbits} applies
also to non-Desarguesian alternative planes and proves the conclusion.

It remains to consider $N=G_2(q)$ on $V=\F_q^6$, with $q$ even and
$q\geq4$. This subgroup is perfect and its commuting field is
$\F_q$, as proved above. Scalar rigidity bounds every alternative
Desarguesian line size by $q$. Any remaining alternative is a
non-Desarguesian affine plane. It cannot have a one-dimensional
flag-transitive group: every subgroup of $A\Gamma L_1(p^d)$ is
soluble, whereas $G$ contains the nonabelian simple subgroup $N$.
The retained BDD list of non-one-dimensional flag-transitive
non-Desarguesian planes consists of the L\"uneburg planes, the nearfield
plane of order nine and the Hering plane of order twenty-seven.
The latter two have an odd number of points, whereas $|V|=q^6$ is
even. Thus only a L\"uneburg plane could remain.

But $G$ is two-transitive on $V$, independently of its incidence
structure. In any invariant linear space this implies two-transitivity
of every line stabilizer on that line: an element sending one ordered
pair of distinct points of the line to another sends their unique
joining line to itself. Hence that line action is primitive, and remains
primitive in the full automorphism group by Lemma~\ref{coarsening}.
This contradicts the retained BDD fact, quoted in source Section~3.2.2,
that the full L\"uneburg-plane group is not locally point-primitive.
This eliminates the last alternative and completes the proof.
\end{proof}

The argument does not assert that two arbitrary invariant affine
structures are nested. Its orbit bounds and classification comparison
replace the unproved nesting sentence and the unnecessary full-group
equality in the printed proof. The classification inputs above remain
charged retained inputs; their proofs have not been reconstructed here.
